\documentclass[UTF8,11pt]{ctexart}
\usepackage[a4paper,margin=22mm]{geometry}
\usepackage{graphicx,enumitem,amssymb,tabularx}
\usepackage[normalem]{ulem}
\usepackage{xeCJKfntef}
\setlength{\parindent}{0pt}
\setlength{\parskip}{.6em}
\sloppy
\begin{document}
\section*{2016年12月英语四级真题(第2套)}
\subsection*{\textbf{2016 年12 月英语四级真题（卷二）}}


\subsection*{\textbf{Part I} \textbf{Writing} \textbf{(30 minutes)}}


Directions: For this part, you are allowed 30 minutes to write an essay. Suppose you have two options upon graduation: one is to work in a state-owned business and the other in a joint venture. You are to make a choice between the two. Write an essay to explain the reasons for your choice. You should write at least\uline{ 120} words but no more than\uline{ 180}words.


\subsection*{\textbf{Part} II \textbf{Listening Comprehension} \textbf{(25 minutes )}}


\subsection*{\textbf{Section A}}


Directions: In this section, you will hear three news reports. At the end of each news report, you will hear two or three questions. Both the news report and the questions will be spoken only once. After you hear a question, you must choose the best answer from the four choices marked A, B, C and D .Then mark the corresponding letter on Answer Sheet I with a single line through the centre.


\textbf{Questions 1 and 2 are based on the news report you have just heard.}


1.


\begin{itemize}[leftmargin=2em,itemsep=.2em,topsep=.4em]\item[A.] To satisfy the curiosity of tourists.
\item[B.] To replace two old stone bridges.
\item[C.] To enable tourists to visit Goat Island.
\item[D.] To improve utility services in the state.
\end{itemize}

2.


\begin{itemize}[leftmargin=2em,itemsep=.2em,topsep=.4em]\item[A.] Countless tree limbs.
\item[B.] A few skeletons.
\item[C.] Lots of wrecked boats and ships.
\item[D.] Millions of coins on the bottom.
\end{itemize}

\textbf{Questions 3 and 4 are based on the news report you have just heard.}


3.


\begin{itemize}[leftmargin=2em,itemsep=.2em,topsep=.4em]\item[A.] It suspended diplomatic relations with Libya.
\item[B.] It urged tourists to leave Tunisia immediately.
\item[C.] It shut down two border crossings with Libya.
\item[D.] It launched a fierce attack against Islamic State.
\end{itemize}

4.


\begin{itemize}[leftmargin=2em,itemsep=.2em,topsep=.4em]\item[A.] Advise Tunisian civilians on how to take safety precautions.
\item[B.] Track down the organization responsible for the terrorist attack.
\item[C.] Train qualified security personnel for the Tunisian government.
\item[D.] Devise a monitoring system on the Tunisian border with Libya.
\end{itemize}

\textbf{Questions 5 to 7 are based on the news report you have just heard.}


5.


\begin{itemize}[leftmargin=2em,itemsep=.2em,topsep=.4em]\item[A.] An environment-friendly battery.
\item[B.] An energy-saving mobile phone.
\item[C.] A plant-powered mobile phone charger.
\item[D.] A device to help plants absorb sunlight.
\end{itemize}

6.


\begin{itemize}[leftmargin=2em,itemsep=.2em,topsep=.4em]\item[A.] While sitting in their school's courtyard.
\item[B.] While playing games on their phones.
\item[C.] While solving a mathematical problem.
\item[D.] While doing a chemical experiment.
\end{itemize}

7.


\begin{itemize}[leftmargin=2em,itemsep=.2em,topsep=.4em]\item[A.] It increases the applications of mobile phones.
\item[B.] It speeds up the process of photosynthesis.
\item[C.] It improves the reception of mobile phones.
\item[D.] It collects the energy released by plants.
\end{itemize}

\subsection*{\textbf{Section B}}


Directions: In this section, you will hear two long conversations. At the end of each conversation, you will hear four questions. Both the conversation and the questions will be spoken only once. After you hear a question, you must choose the best answer from the four choices marked A, B, C and D. Then mark the corresponding letter on Answer Sheet 1with a single line through the centre.


\textbf{Questions 8 to 11 are based on the conversation you have just heard.}


8.


\begin{itemize}[leftmargin=2em,itemsep=.2em,topsep=.4em]\item[A.] He visited the workshops in the Grimsby plant.
\item[B.] He called the woman and left her a message.
\item[C.] He used stand-ins as replacements on all lines.
\item[D.] He asked a technician to fix the broken production line.
\end{itemize}

9.


\begin{itemize}[leftmargin=2em,itemsep=.2em,topsep=.4em]\item[A.] It is the most modern production line.
\item[B.] It assembles super-intelligent robots.
\item[C.] It has stopped working completely.
\item[D.] It is going to be upgraded soon.
\end{itemize}

10.


\begin{itemize}[leftmargin=2em,itemsep=.2em,topsep=.4em]\item[A.] To seek her permission.
\item[B.] To place an order for robots.
\item[C.] To request her to return at once.
\item[D.] To ask for Tom's phone number.
\end{itemize}

\noindent\begin{tabularx}{\linewidth}{@{}lXXXX@{}}11. & \textbf{A.} She is on duty. & \textbf{B.} She is having her day off. & \textbf{C.} She is on sick leave. & \textbf{D.} She is abroad on business.\\\end{tabularx}\par

\textbf{Questions 12 to 15 are based on the conversation you have just heard.}


12.


\begin{itemize}[leftmargin=2em,itemsep=.2em,topsep=.4em]\item[A.] He saved a baby boy's life.
\item[B.] He wanted to be a superhero.
\item[C.] He prevented a train crash.
\item[D.] He was a witness to an accident.
\end{itemize}

13.


\begin{itemize}[leftmargin=2em,itemsep=.2em,topsep=.4em]\item[A.] He has a 9-month-old boy.
\item[B.] He is currently unemployed.
\item[C.] He enjoys the interview.
\item[D.] He commutes by subway.
\end{itemize}

\noindent\begin{tabularx}{\linewidth}{@{}lXXXX@{}}14. & \textbf{A.} A rock on the tracks. & \textbf{B.} A misplaced pushchair. & \textbf{C.} A strong wind. & \textbf{D.} A speeding car.\\\end{tabularx}\par

15.


\begin{itemize}[leftmargin=2em,itemsep=.2em,topsep=.4em]\item[A.] She stood motionless in shock.
\item[B.] She cried bitterly.
\item[C.] She called the police at once.
\item[D.] She shouted for help.
\end{itemize}

\subsection*{\textbf{Section C}}


Directions: In this section, you will hear three passages. At the end of each passage, you will hear three or four questions. Both the passage and the questions will be spoken only once. After you hear a question, you must choose the best answer from the four choices marked A, B, C and D. Then mark the corresponding letter on Answer Sheet I with a single line through the centre.


\textbf{Questions 16 to 18 are based on the passage you have just heard.}


16.


\begin{itemize}[leftmargin=2em,itemsep=.2em,topsep=.4em]\item[A.] She inherited her family ice-cream business in Billings.
\item[B.] She loved the ice-cream business more than teaching primary school.
\item[C.] She started an ice-cream business to finance her daughter's education.
\item[D.] She wanted to have an ice-cream truck when she was a little girl.
\end{itemize}

\noindent\begin{tabularx}{\linewidth}{@{}lXXXX@{}}17. & \textbf{A.} To preserve a tradition. & \textbf{B.} To amuse her daughter. & \textbf{C.} To help local education. & \textbf{D.} To make some extra money.\\\end{tabularx}\par

18.


\begin{itemize}[leftmargin=2em,itemsep=.2em,topsep=.4em]\item[A.] To raise money for business expansion.
\item[B.] To make her truck attractive to children.
\end{itemize}

\begin{itemize}[leftmargin=2em,itemsep=.2em,topsep=.4em]\item[C.] To allow poor kids to have ice-cream too.
\item[D.] To teach kids the value of mutual support.
\end{itemize}

\textbf{Questions 19 to 21 are based on the passage you have just heard.}


19.


\begin{itemize}[leftmargin=2em,itemsep=.2em,topsep=.4em]\item[A.] The reasons for imposing taxes.
\item[B.] The various services money can buy.
\item[C.] The various burdens on ordinary citizens.
\item[D.] The function of money in the modem world.
\end{itemize}

20.


\begin{itemize}[leftmargin=2em,itemsep=.2em,topsep=.4em]\item[A.] Educating and training citizens.
\item[B.] Improving public transportation.
\item[C.] Protecting people's life and property.
\item[D.] Building hospitals and public libraries.
\end{itemize}

21.


\begin{itemize}[leftmargin=2em,itemsep=.2em,topsep=.4em]\item[A.] By asking for donations.
\item[B.] By selling public lands.
\item[C.] By selling government bonds.
\item[D.] By exploiting natural resources.
\end{itemize}

\textbf{Questions 22 to 25 are based on the passage you have just heard.}


22.


\begin{itemize}[leftmargin=2em,itemsep=.2em,topsep=.4em]\item[A.] It is located at the center of the European continent.
\item[B.] It relies on tourism as its chief source of revenues.
\item[C.] It contains less than a square mile of land.
\item[D.] It is surrounded by France on three sides.
\end{itemize}

23.


\begin{itemize}[leftmargin=2em,itemsep=.2em,topsep=.4em]\item[A.] Its beauty is frequently mentioned in American media.
\item[B.] Its ruler Prince Rainier married an American actress.
\item[C.] It is where many American movies are shot.
\item[D.] It is a favorite place Americans like to visit.
\end{itemize}

\noindent\begin{tabularx}{\linewidth}{@{}lXXXX@{}}24. & \textbf{A.} Tobacco. & \textbf{B.} Potatoes. & \textbf{C.} Machinery. & \textbf{D.} Clothing.\\\end{tabularx}\par

25.


\begin{itemize}[leftmargin=2em,itemsep=.2em,topsep=.4em]\item[A.] European history.
\item[B.] European geography.
\item[C.] Small countries in Europe.
\item[D.] Tourist attractions in Europe.
\end{itemize}

\subsection*{\textbf{Part} III \textbf{Reading Comprehension} \textbf{(40 minutes)}}


\subsection*{\textbf{Section A}}


Directions: In this section, there is a passage with ten blanks. You are required to select one word for each blank from a list of choices given in a word bank following the passage. Read the passage through carefully before making your choices. Each choice in the bank is identified by a letter. Please mark the corresponding letter for each item on Answer Sheet 2 with a single line through the centre. You may not use any of the words in the bank more than once.


\textbf{Questions 26 to 35 are based on the following passage.}


The ocean is heating up. That's the conclusion of a new study that finds that Earth's oceans now\uline{ 26 }heat at twice the rate they did 18 years ago. Around half of ocean heat intake since 1865 has taken place since 1997, researchers report online in Nature Climate Change.


Warming waters are known to\uline{ 27 }to coral bleaching (珊瑚白化) and they take up more space than cooler waters, raising sea\uline{ 28 }While the top of the ocean is well studied, its depths are more difficult to\uline{ 29 }The researchers gathered 150 years of ocean temperature data in order to get abetter\uline{ 30 }of heat absorption from surface to seabed. They gathered together temperature readings collected by everything from a 19th century\uline{ 31} of British naval ships to modem automated ocean probes. The extensive data sources,\uline{ 32 }with computer simulations ( 计算机模拟), created a timeline of ocean temperature changes, including cooling from volcanic outbreaks and warming from fossil fuel\uline{ 33}


About 35 percent of the heat taken in by the oceans during the industrial era now resides at a\uline{ 34 }of more than 700 meters, the researchers found. They say they're\uline{ 35 }whether the deep-sea warming canceled out warming at the sea's surface.


\begin{itemize}[leftmargin=2em,itemsep=.2em,topsep=.4em]\item[A.] absorb
\item[B.] Combined
\item[C.] Contribute
\item[D.] depth
\item[E.] emissions
\item[F.] excursion
\item[G.] explore
\item[H.] floor
\item[I.] heights
\item[J.] indifferent
\item[K.] levels
\item[L.] mixed
\item[M.] picture
\item[N.] unsure
\item[O.] voyage
\end{itemize}

\subsection*{\textbf{Section B}}


Directions: In this section, you are going to read a passage with ten statements attached to it. Each statement contains information given in one of the paragraphs. Identify the paragraph from which the information is derived. You may choose a paragraph more than once. Each paragraph is marked with a letter. Answer the questions by marking the corresponding letter on Answer Sheet 2. \textbf{The Secret to Raising Smart Kids}


\begin{itemize}[leftmargin=2em,itemsep=.2em,topsep=.4em]\item[A.] I first began to investigate the basis of human motivation--and how people persevere after setbacks--as
\end{itemize}

a psychology graduate student at Yale University in the 1960s. Animal experiments by psychologists at the University of Pennsylvania had shown that after repeated failures, most animals conclude that a situation is hopeless and beyond their control. After such an experience an animal often remains passive even when it can effect change--a state they called learned helplessness.


\begin{itemize}[leftmargin=2em,itemsep=.2em,topsep=.4em]\item[B.] People can learn to be helpless, too. Why do some students give up when they encounter difficulty,
\end{itemize}

whereas others who are no more skilled continue to strive and learn? One answer, I soon discovered, lay in people's beliefs about why they had failed.


\begin{itemize}[leftmargin=2em,itemsep=.2em,topsep=.4em]\item[C.] In particular, attributing poor performance to a lack of ability depresses motivation more than does the
\end{itemize}

belief that lack of effort is to blame. When I told a group of school children who displayed helpless behavior that a lack of effort led to their mistakes in math, they learned to keep trying when the problems got tough. Another group of helpless children who were simply rewarded for their success on easier problems did not improve their ability to solve hard math problems. These experiments indicated that a focus on effort can help resolve helplessness and generate success.


\begin{itemize}[leftmargin=2em,itemsep=.2em,topsep=.4em]\item[D.] Later, I developed a broader theory of what separates the two general classes of learners—helpless
\end{itemize}

versus mastery-oriented. I realized these different types of students not only explain their failures differently, but they also hold different "theories" of intelligence. The helpless ones believe intelligence is a fixed characteristic: you have only a certain amount, and that's that. I call this a "fixed mind-set (思维模式). " Mistakes crack their self-confidence because they attribute errors to a lack of ability, which they feel powerless to change. They avoid challenges because challenges make mistakes more likely. The mastery-oriented children, on the other hand, think intelligence is not fixed and can be developed through education and hard work. Such children believe challenges are energizing rather than intimidating (令人生畏); they offer opportunities to learn. Students with such a growth mind-set were destined (注定) for greater academic success and were quite likely to outperform their counterparts.


\begin{itemize}[leftmargin=2em,itemsep=.2em,topsep=.4em]\item[E.] We validated these expectations in a study in which two other psychologists and I monitored
\end{itemize}

373students for two years during the transition to junior high school, when the work gets more difficult and the grading more strict, to determine how their mind-sets might affect their math grades. At the beginning of seventh grade, we assessed the students' mind-sets by asking them to agree or disagree with statements such as "Your intelligence is something very basic about you that you can't really change. " We then assessed their beliefs about other aspects of learning and looked to see what happened to their grades.


\begin{itemize}[leftmargin=2em,itemsep=.2em,topsep=.4em]\item[F.] As predicted, the students with a growth mind-set felt that learning was a more important goal than
\end{itemize}

getting good grades. In addition, they held hard work in high regard. They understood that even geniuses have to work hard. Confronted by a setback such as a disappointing test grade, students with a growth mind-set said they would study harder or try a different strategy. The students who held a fixed mind-set, however, were concerned about looking smart with less regard for learning. They had negative views of effort, believing that having to work hard was a sign of low ability. They thought that a person with talent or intelligence did not need to work hard to do well. Attributing a bad grade to their own lack of ability, those with a fixed mind-set said that they would study less in the future, try never to take that subject again and consider cheating on future tests.


\begin{itemize}[leftmargin=2em,itemsep=.2em,topsep=.4em]\item[G.] Such different outlooks had a dramatic impact on performance. At the start of junior high, the math
\end{itemize}

achievement test scores of the students with a growth mind-set were comparable to those of students who displayed a fixed mind-set. But as the work became more difficult, the students with a growth mind-set showed greater persistence. As a result, their math grades overtook those of the other students by the end of the first semester--and the gap between the two groups continued to widen during the two years we followed them.


\begin{itemize}[leftmargin=2em,itemsep=.2em,topsep=.4em]\item[H.] A fixed mind-set can also hinder communication and progress in the workplace and discourage or
\end{itemize}

ignore constructive criticism and advice. Research shows that managers who have a fixed mind-set are less likely to seek or welcome feedback from their employees than are managers with a growth mind-set.


\begin{itemize}[leftmargin=2em,itemsep=.2em,topsep=.4em]\item[I.] How do we transmit a growth mind-set to our children? One way is by telling stories about achievements
\end{itemize}

that result from hard work. For instance, talking about mathematical geniuses who were more or less born that way puts students in a fixed mind-set, but descriptions of great mathematicians who fell in love with math and developed amazing skills produce a growth mind-set.


\begin{itemize}[leftmargin=2em,itemsep=.2em,topsep=.4em]\item[J.] In addition, parents and teachers can help children by providing explicit instruction regarding the mind
\end{itemize}

as a learning machine. I designed an eight-session workshop for 91 students whose math grades were declining in their first year of junior high. Forty-eight of the students received instruction in study skills only, whereas the others attended a combination of study skills sessions and classes in which they learned about the growth mind-set and how to apply it to schoolwork. In the growth mind-set classes, students read and discussed an article entitled "You Can Grow Your Brain. " They were taught that the brain is like a muscle that gets stronger with use and that learning prompts the brain to grow new connections. From such instruction, many students began to see themselves as agents of their own brain development. Despite being unaware that there were two types of instruction, teachers reported significant motivational changes in 27\% of the children in the growth mind-set workshop as compared with only 9\% of students in the control group.


\begin{itemize}[leftmargin=2em,itemsep=.2em,topsep=.4em]\item[K.] Research is converging (汇聚) on the conclusion that great accomplishment and even genius is typically
\end{itemize}

the result of years of passion and dedication and not something that flows naturally from a gift.


36. The author's experiment shows that students with a fixed mind-set believe having to work hard is an


indication of low ability.


37. Focusing on effort is effective in helping children overcome frustration and achieve success.


38. We can cultivate a growth mind-set in children by telling success stories that emphasize hard work and


love of learning.


39. Students' belief about the cause of their failure explains their attitude toward setbacks.


40. In the author's experiment, students with a growth mind-set showed greater perseverance in solving


difficult math problems.


41. The author conducted an experiment to find out about the influence of students' mind-sets on math


learning.


42. After failing again and again, most animals give up hope.


43. Informing students about the brain as a learning machine is a good strategy to enhance their motivation


for learning.


44. People with a fixed mind-set believe that one's intelligence is unchangeable.


45. In the workplace, feedback may not be so welcome to managers with a fixed mind-set.


\subsection*{\textbf{Section C}}


Directions: There are 2 passages in this section. Each passage is followed by some questions or unfinished statements. For each of them there are four choices marked A., B., C. And D . You should decide on the best choice and mark the corresponding letter on Answer Sheet 2 with a single line through the centre.


\subsection*{\textbf{Passage One}}


\textbf{Questions 46 to 50 are based on the following passage.}


"Sugar, alcohol and tobacco," economist Adam Smith once wrote," are commodities which are nowhere necessaries of life, which have become objects of almost universal consumption, and which are, therefore, extremely popular subjects of taxation. "


Two and a haft centuries on, most countries impose some sort of tax on alcohol and tobacco. With surging obesity levels putting increasing strain on public health systems, governments around the world have begun to toy with the idea of taxing sugar as well.


Whether such taxes work is a matter of debate. A preliminary review of Mexico's taxation found a fall in purchases of taxed drinks as well as a rise in sales of untaxed and healthier drinks. By contrast, a Danish tax on foods high in fats was abandoned a year after its introduction, amid claims that consumers were avoiding it by crossing the border to Germany to satisfy their desire for cheaper, fattier fare.


The food industry has, in general, been firmly opposed to such direct government action. Nonetheless, the renewed focus on waistlines means that industry groups are under pressure to demonstrate their products are healthy as well as tasty.


Over the past three decades, the industry has made some efforts to improve the quality of its offerings. For example, some drink manufacturers have cut the amount of sugar in their beverages.


Many of the reductions over the past 30 years have been achieved either by reducing the amount of sugar, salt or fat in a product, or by finding an alternative ingredient. More recently, however, some companies have been investing money in a more ambitious undertaking: learning how to adjust the fundamental make-up of the food they sell. For example, having salt on the outside, but none on the inside, reduces the salt content without changing the taste.


While reformulating recipes ( 配方) is one way to improve public health, it should be part of a multi-sided approach. The key is to remember that there is not just one solution. To deal with obesity, mixture of approaches--including reformulation, taxation and adjusting portion sizes--will be needed. There is no silver bullet.


46. What did Adam Smith say about sugar, alcohol and .tobacco?


\begin{itemize}[leftmargin=2em,itemsep=.2em,topsep=.4em]\item[A.] They were profitable to manufacture.
\item[B.] They were in ever-increasing demand.
\item[C.] They were subject to taxation almost everywhere.
\item[D.] They were no longer considered necessities of life.
\end{itemize}

47. Why have many countries started to consider taxing sugar?


\begin{itemize}[leftmargin=2em,itemsep=.2em,topsep=.4em]\item[A.] They are under growing pressure to balance their national budgets.
\item[B.] They find it ever harder to cope with sugar-induced health problems.
\item[C.] The practice of taxing alcohol and tobacco has proved both popular and profitable.
\item[D.] The sugar industry is overtaking alcohol and tobacco business in generating profits.
\end{itemize}

48. What do we learn about Danish taxation on fat-rich foods?


\begin{itemize}[leftmargin=2em,itemsep=.2em,topsep=.4em]\item[A.] It did not work out as well as was expected.
\item[B.] It gave rise to a lot of problems on the border.
\item[C.] It could not succeed without German cooperation.
\item[D.] It met with firm opposition from the food industry.
\end{itemize}

49. What is the more recent effort by food companies to make foods and drinks both healthy and tasty?


\begin{itemize}[leftmargin=2em,itemsep=.2em,topsep=.4em]\item[A.] Replacing sugar or salt with alternative ingredients.
\end{itemize}

\begin{itemize}[leftmargin=2em,itemsep=.2em,topsep=.4em]\item[B.] Setting a limit on the amount of sugar or salt in their products.
\item[C.] Investing in research to find ways to adapt to consumers' needs.
\item[D.] Adjusting the physical composition of their products.
\end{itemize}

50. What does the author mean by saying, at the end of the passage, "There is no silver bullet" ( Line 4,


Para.7)?


\begin{itemize}[leftmargin=2em,itemsep=.2em,topsep=.4em]\item[A.] There is no single easy quick solution to the problem.
\item[B.] There is no hope of success without public cooperation.
\item[C.] There is no hurry in finding ways to solve the obesity problem.
\item[D.] There is no effective way to reduce people's sugar consumption.
\end{itemize}

\subsection*{\textbf{Passage Two}}


\textbf{Questions 51 to 55 are based on the following passage.}


You may have heard some of the fashion industry horror stories: models eating tissues or cotton balls to hold off hunger, and models collapsing from hunger-induced heart attacks just seconds after they step off the runway.


Excessively skinny models have been a point of controversy for decades, and two researchers say a model's body mass should be a workplace health and safety issue. In an editorial released Monday in the American Journal of Public Health, Katherine Record and Bryn Austin made their case for government regulation of the fashion industry.


The average international runway model has a body mass index (BMI) under 16--low enough to indicate starvation by the World Health Organization's standard. And Record and Austin are worried not just about the models themselves, but about the vast number of girls and women their images influence.


"Especially girls and teens," says Record. "Seventy percent of girls aged 10 to 18 report that they define perfect body image based on what they see in magazines. " That's especially worrying, she says, given that anorexia (厌食症) results in more deaths than does any other mental illness, according to the National Institute of Mental Health.


It's commonly known that certain diseases are linked with occupations like lung disease in coalminers. Professional fashion models are particularly vulnerable to eating disorders resulting from occupational demands to maintain extreme thinness.


Record's suggestion is to prohibit agents from hiring models with a BMI below 18.


In April, France passed a law setting lower limits for a model's weight. Agents and fashion houses who hire models with a BMI under 18 could pay \$ 82,000 in fines and spend up to 6 months in jail. Regulating the fashion industry in the United States won't be easy, Record says. But with the new rules in France, U.S. support could make a difference. "A designer can't survive without participating in Paris Fashion Week", she says, adding," Our argument is that the same would be true of New York Fashion Week. "


51. What do Record and Austin say about fashion models' body mass?


\begin{itemize}[leftmargin=2em,itemsep=.2em,topsep=.4em]\item[A.] It has caused needless controversy.
\item[B.] It is but a matter of personal taste.
\item[C.] It is the focus of the modeling business.
\item[D.] It affects models' health and safety.
\end{itemize}

52. What are Record and Austin advocating in the Monday editorial?


\begin{itemize}[leftmargin=2em,itemsep=.2em,topsep=.4em]\item[A.] A change in the public's view of female beauty.
\end{itemize}

\begin{itemize}[leftmargin=2em,itemsep=.2em,topsep=.4em]\item[B.] Government legislation about models' weight.
\item[C.] Elimination of forced weight loss by models.
\item[D.] Prohibition of models eating non-food stuff.
\end{itemize}

53. Why are Record and Austin especially worried about the low body mass index of models?


\begin{itemize}[leftmargin=2em,itemsep=.2em,topsep=.4em]\item[A.] It contributes to many mental illnesses.
\item[B.] It defines the future of the fashion industry.
\item[C.] It has great influence on numerous girls and women.
\item[D.] It keeps many otherwise qualified women off the runway.
\end{itemize}

54. What do we learn about France's fashion industry?


\begin{itemize}[leftmargin=2em,itemsep=.2em,topsep=.4em]\item[A.] It has difficulty hiring models.
\item[B.] It has now a new law to follow.
\item[C.] It allows girls under 18 on the runway.
\item[D.] It has overtaken that of the United States.
\end{itemize}

55. What does Record expect of New York Fashion Week?


\begin{itemize}[leftmargin=2em,itemsep=.2em,topsep=.4em]\item[A.] It will create a completely new set of rules.
\item[B.] It will do better than Paris Fashion Week.
\item[C.] It will differ from Paris Fashion Week.
\item[D.] It will have models with a higher BMI.
\end{itemize}

\textbf{Part} IV \textbf{Translation} \textbf{( 30 minutes )} Directions: For this part, you are allowed 30 minutes to translate a passage from Chinese into English. You should write your answer on Answer Sheet 2.


随着中国的改革开放，如今很多年轻人都喜欢举行西式婚礼。新娘在婚礼上穿着白色婚纱，因为白色被认为是纯洁的象征。然而，在中国传统文化中，白色经常是葬礼上使用的颜色。因此务必记住，白花一定不要用作祝人康复的礼物，尤其不要送给老年人或危重病人。同样，礼金也不能装在白色的信封里，而要装在红色的信封里。

\end{document}